%********************************************************************
% Overall conclusion and future work
%********************************************************************
\chapter{Overall Conclusion and Future Work}

\section{Overall Conclusion}
The promises of scaling laws, where scaling up model size gives rise to new capabilities, 
are colliding with a harsh reality: training these large models and running inference
are costly. The central research question is how to scale language models efficiently,
both for training and inference. It is fortunate that pre-trained LLMs exhibit
redundancy at different levels that can be exploited to improve their efficiency
during post-training or pretraining. One example of
such redundancy is the low-rank structure of pretrained Transformer activations.
Our contributions exploit this observation to improve LLM efficiency in two main
directions: \textbf{model compression} using activations and \textbf{activation sparsity}.

For post-training \textbf{model compression}, we propose \textbf{Lillama}, a method that locally distills the
activations of Transformer layers into low-rank weights. This local optimization makes
the approach more efficient than prior work: it enables compressing, on a single GPU,
LLMs that would not otherwise fit on one. To achieve speedup, the method does not require
specialized GPU kernels, unlike weight-sparsity approaches such as SparseGPT~\citep{frantar2023sparsegpt} or Wanda~\citep{sun2024wanda}.
The approach is also sample-efficient, requiring only 13M tokens, compared to methods
like Sheared-Llama \citep{xia2024sheared}, or Minitron \citep{muralidharan2024compact,
sreenivas2024minitron}, which use over 50B tokens for compression. Sample efficiency is an important property in post-training compression, as it
directly affects the fairness of the approach---for instance, its viability for
under-represented languages. This is why we specifically studied Whisper compression in
a low-resource setting, using Bambara as a case study. In this work, \textbf{BaldWhisper}, we
again distill input/output embedding activations into low-rank weights, but push
compression further by reducing the number of layers. To preserve the method's sample
efficiency, we do this by merging layers rather than removing them,
to stay close to the original model.
The resulting model is twice as fast while maintaining most of the performance.

For \textbf{activation sparsity}, we first study the eigenstructure of \textbf{Matryoshka
embedding LLMs} in multilingual and multimodal representation learning. We show
that Matryoshka embeddings, viewed as a form of structured activation sparsity,
leave room for more effective use of the representation space: most of the
activation energy is concentrated in relatively few principal components.
The fixed slicing operator applied to the embedding vectors
does not preserve all the information needed for downstream tasks while
introducing redundancy. We make the same observation for dense MLPs in
autoregressive LLMs: their intermediate activations occupy a subspace
substantially smaller than the allocated dimension. Sparse MoEs, however,
project each token into smaller subspaces of the available intermediate space,
and we show that these smaller subspaces are used more fully.
These observations motivate an approach where the MLP has a large intermediate
dimension but each token adaptively selects only the intermediate coordinates it needs
rather than activating the whole available space.
We introduce \textbf{ZeroNeuron (ZoN)}, an input-dependent
gating mechanism that learns which coordinates to activate while constraining the average active dimension.
Compared to fine-grained sparse mixture-of-experts and sparse memory layers,
ZoN gets rid of TopK and predicts sparsity pointwise, which
makes it not only more compute-adaptive but also more communication-efficient
in distributed training, compared to TopK-based sparse mixture-of-experts and
sparse memory layers that require an all-to-all token dispatch and all-to-all return.
With approximately 600M active parameters per
token, the 1.2B ZoN model matches the aggregate performance of its dense
counterpart and outperforms a dense model of a comparable number of active parameters.
These results show that model capacity can be preserved while reducing
the parameters activated for each input, and that such sparsity
can be exploited by GPU kernels to accelerate training and inference.

\section{Future Work}
We believe in sparsity, and we expect the most performant machine learning systems
of the future to be sparsely activated. The remaining question is how
this sparsity should be designed. We do not believe that TopK-based
mixtures of experts or sparse memory layers are the right long-term
answer: there is no reason to impose the same activation budget on
every token, regardless of the task and its difficulty. Moreover, under
current hardware constraints, distributed TopK routing incurs substantial
cross-GPU expert communication. This is one of the biggest challenges in
scaling fine-grained sparsity, and approaches such as LatentMoE
~\citep{elango2026latentmoeoptimalaccuracyflop} are trying to address
it. Maybe this will no longer be a problem in a few years, if the next generation
of hardware offers individual GPUs with terabytes of VRAM and the capacity to train
LLMs with trillions of parameters on a single device.
But beyond the engineering issues, TopK-based
routing also faces the stability problems discussed earlier, including
dead experts and expert collapse. Rather than putting patches everywhere
to keep TopK-based sparsity working, we believe it can be beneficial to fundamentally
revisit how sparsity is learned. We believe that such approaches must be adaptive, communication-efficient,
and stable by design. ZoN is our first step in this direction. We plan
to investigate its scalability through larger models and longer training
runs, while developing GPU kernels that translate its learned sparsity
into actual computational savings.